\section{维数正规化下的能动量张量}
\label{sec:2}
在规范理论中，能动量张量的描述~\cite{Freedman:1974gs,Joglekar:1975jm}在维数正规化下
尤为简单。\footnote{文献~\cite{Collins:1984xc}是对维数正规化的一份非常好的阐述。}
这是因为该正规化明显保持了（矢量）规范对称性与平移不变性。因此本节基于维数正规化简要
复述与能动量张量有关的基本事实。

所考虑系统在 $D$ 维欧氏空间中的作用量为
\begin{equation}
   S=\frac{1}{4g_0^2}\int\mathrm{d}^Dx\,F_{\mu\nu}^a(x)F_{\mu\nu}^a(x)
   +\int\mathrm{d}^Dx\,\Bar{\psi}(x)(\Slash{D}+m_0)\psi(x),
\label{eq:(2.1)}
\end{equation}
其中 $g_0$ 与~$m_0$ 分别是裸规范耦合与质量参数。场强定义为
\begin{equation}
   F_{\mu\nu}(x)
   =\partial_\mu A_\nu(x)-\partial_\nu A_\mu(x)+[A_\mu(x),A_\nu(x)],
\label{eq:(2.2)}
\end{equation}
其中 $A_\mu(x)=A_\mu^a(x)T^a$ 及~$F_{\mu\nu}(x)=F_{\mu\nu}^a(x)T^a$；作用在费米子上的
协变导数为
\begin{equation}
   D_\mu=\partial_\mu+A_\mu.
\label{eq:(2.3)}
\end{equation}
这里以及下文中，对 $N_{\mathrm{f}}$ 个费米子味道的求和一律略去不写。我们的 $\gamma$
矩阵取为厄米的，旋量指标的迹对任何~$D$ 都取 $\tr 1=4$。我们还取
\begin{equation}
   D=4-2\epsilon.
\label{eq:(2.4)}
\end{equation}

在维数正规则化的假定下，与平移不变性相联系的 Ward--Takahashi 关系可以直截了当地推导。
我们考虑对泛函积分中积分变量的下列无穷小变分：
\begin{equation}
   \delta A_\mu(x)=\xi_\nu(x)F_{\nu\mu}(x),\qquad
   \delta\psi(x)=\xi_\mu(x)D_\mu\psi(x).
\label{eq:(2.5)}
\end{equation}
于是，由于作用量的变化为\footnote{为了使微扰论良定义，这里隐含地假定存在规范固定项与
Faddeev--Popov 鬼场。但由于它们不以显式方式出现在规范不变算符的关联函数中，
下文中我们将略去这些要素。}
\begin{equation}
   \delta S=-\int\mathrm{d}^Dx\,\xi_\nu(x)\partial_\mu T_{\mu\nu}(x),
\label{eq:(2.6)}
\end{equation}
其中能动量张量~$T_{\mu\nu}(x)$ 定义为
\begin{align}
   T_{\mu\nu}(x)
   &\equiv\frac{1}{g_0^2}\left[
   F_{\mu\rho}^a(x)F_{\nu\rho}^a(x)
   -\frac{1}{4}\delta_{\mu\nu}F_{\rho\sigma}^a(x)F_{\rho\sigma}^a(x)
   \right]
\notag\\
   &\qquad{}
   +\frac{1}{4}
   \Bar{\psi}(x)\left(\gamma_\mu\overleftrightarrow{D}_\nu
   +\gamma_\nu\overleftrightarrow{D}_\mu\right)\psi(x)
   -\delta_{\mu\nu}\Bar{\psi}(x)
   \left(\frac{1}{2}\overleftrightarrow{\Slash{D}}
   +m_0\right)\psi(x),
\label{eq:(2.7)}
\end{align}
且
\begin{equation}
   \overleftrightarrow{D}_\mu\equiv D_\mu-\overleftarrow{D}_\mu,\qquad
   \overleftarrow{D}_\mu\equiv\overleftarrow{\partial}_\mu-A_\mu,
\label{eq:(2.8)}
\end{equation}
我们有
\begin{equation}
   \left\langle\mathcal{O}_{\text{out}}
   \int_{\mathcal{D}}
   \mathrm{d}^Dx\,\partial_\mu T_{\mu\nu}(x)\,\mathcal{O}_{\text{in}}\right\rangle
   =-\left\langle\mathcal{O}_{\text{out}}\,\partial_\nu\mathcal{O}_{\text{in}}
   \right\rangle,
\label{eq:(2.9)}
\end{equation}
其中 $\mathcal{O}_{\text{in}}$（$\mathcal{O}_{\text{out}}$）是局域在有限积分
区域~$\mathcal{D}$ 内部（外部）的一组规范不变算符。该关系表明能动量张量产生无穷小平移，
同时裸量~\eqref{eq:(2.7)} 不接受相乘重正化。于是，通过减去其（可能发散的）真空期望值来
定义重正化的有限能动量张量：
\begin{equation}
   \left\{T_{\mu\nu}\right\}_R(x)
   \equiv T_{\mu\nu}(x)-\left\langle T_{\mu\nu}(x)\right\rangle.
\label{eq:(2.10)}
\end{equation}

能动量张量的一个基本性质是迹
反常~\cite{Crewther:1972kn,Chanowitz:1972vd,Adler:1976zt,Nielsen:1977sy,%
Collins:1976yq}。推导它的一个简单办法~\cite{Fujikawa:1980rc,Fujikawa:2004cx}是在式%
~\eqref{eq:(2.5)} 中取 $\xi_\mu(x)\propto x_\mu$ 并把所得关系与重正化群方程比较。
经过一些考虑，得到
\begin{equation}
   \delta_{\mu\nu}
   \left\{T_{\mu\nu}\right\}_R(x)
   =-\frac{\beta}{2g^3}\left\{F_{\rho\sigma}^aF_{\rho\sigma}^a\right\}_R(x)
   -(1+\gamma_m)m\left\{\Bar\psi\psi\right\}_R(x),
\label{eq:(2.11)}
\end{equation}
其中我们假定右边的重正化算符按
$\text{MS}$ 方案定义~\cite{Collins:1984xc}。\footnote{$\text{MS}$ 方案下单圈重正化的
内容总结于附录~\ref{sec:A}。}在本文中我们总是假定重正化算符中已减去真空期望值。
在式~\eqref{eq:(2.11)} 中，重正化群函数 $\beta$ 和~$\gamma_m$ 定义为
\begin{align}
   \beta&\equiv
   \left(\mu\frac{\partial}{\partial\mu}\right)_0g
   =-\frac{1}{2}g\left(\mu\frac{\partial}{\partial\mu}\right)_0\ln Z,
\label{eq:(2.12)}\\
   \gamma_m&\equiv
   -\left(\mu\frac{\partial}{\partial\mu}\right)_0\ln m
   =-\beta\frac{\partial}{\partial g}\ln Z_m,
\label{eq:(2.13)}
\end{align}
其中 $g$ 与~$m$ 分别是重正化规范耦合与重正化质量，而对重正化标度~$\mu$
的求导是在所有裸量保持固定的条件下进行的。重正化常数定义为
\begin{equation}
   g_0^2=\mu^{2\epsilon}g^2Z,\qquad m_0=mZ_m^{-1}.
\label{eq:(2.14)}
\end{equation}
这些重正化函数微扰展开的前几项为
\begin{align}
   \beta=-b_0g^3-b_1g^5+O(g^7),\qquad\gamma_m=d_0g^2+d_1g^4+O(g^6),
\label{eq:(2.15)}
\end{align}
其中~\cite{Caswell:1974gg,Jones:1974mm}
\begin{align}
   b_0&=\frac{1}{(4\pi)^2}
   \left[\frac{11}{3}C_2(G)-\frac{4}{3}T(R)N_{\mathrm{f}}\right],
\label{eq:(2.16)}\\
   b_1&=\frac{1}{(4\pi)^4}
   \left\{
   \frac{34}{3}C_2(G)^2-\left[4C_2(R)+\frac{20}{3}C_2(G)\right]T(R)N_{\mathrm{f}}
   \right\},
\label{eq:(2.17)}
\end{align}
以及~\cite{Tarrach:1980up,Nachtmann:1981zg}
\begin{align}
   d_0&=\frac{1}{(4\pi)^2}6C_2(R),
\label{eq:(2.18)}\\
   d_1&=\frac{1}{(4\pi)^4}
   \left\{
   3C_2(R)^2
   +\left[\frac{97}{3}C_2(G)-\frac{20}{3}T(R)N_{\mathrm{f}}\right]C_2(R)
   \right\}.
\label{eq:(2.19)}
\end{align}
